\section{Related Work}
\label{sec:related}

\paragraph{LLM serving systems.}
Orca~\cite{yu2022orca} introduced iteration-level scheduling, and
\vllm~\cite{kwon2023pagedattention} added paged KV-cache management; together they define
the continuous-batching design we study.
Sarathi-Serve~\cite{agrawal2024sarathi} splits prefills into chunks to protect decode
latency, and DistServe~\cite{zhong2024distserve} places prefill and decode on different
GPUs.
FlashAttention~\cite{dao2022flashattention,dao2024flashattention2} reduces the memory
traffic of attention kernels.
These systems improve scheduling and kernels; we instead measure and model how one of
them scales with the GPU interconnect.

\paragraph{Model parallelism.}
Megatron-LM~\cite{shoeybi2019megatron} introduced the TP partitioning used by \vllm, and
Narayanan et al.~\cite{narayanan2021efficient} combined tensor, pipeline and data
parallelism for training at scale.
Their evaluation targets training on NVLink and InfiniBand clusters, where messages are
large; inference decode exchanges much smaller messages.

\paragraph{Performance modeling.}
The roofline model~\cite{williams2009roofline} bounds performance by compute and memory
bandwidth, and latency--bandwidth models of collectives~\cite{hockney1994communication,
thakur2005optimization} estimate the cost of all-reduce.
Pope et al.~\cite{pope2023efficiently} combine both ideas to choose partitioning layouts
for Transformer inference on TPU~v4.
Our model is simpler and is calibrated and validated on GPUs, against an open-source
serving engine, on two very different interconnects.
